\section{Experimental Design}
\label{sec:experiments}

\subsection{Stage 0: correctness gates}

Before training image classifiers, we use a low-dimensional strongly convex
problem for which the exact hypergradient is available. We compare (i) the
analytic hypergradient, (ii) finite differences, (iii) fully unrolled
autodifferentiation, and (iv) Eq.~\eqref{eq:loss-difference-gradient}. The
sign must agree and pairwise cosine similarity must exceed $0.99$. We also
test that projection preserves $0\leq\omega_i\leq1$ and
$\sum_i\omega_i=B$, that the uniform state equals $B/n$, and that tie breaking
is deterministic.

On MNIST with 10\% label noise, we then hold $\vomega$ fixed and increase the
number of inner updates. We record both inner optimality residuals, score
cosine, Top-$B$ Jaccard similarity, and agreement with an independently
retrained solution. The main study starts only after the selected stopping
rule produces a stable score. Validation data determine stopping and
hyperparameters; the test set is never used for either.

\subsection{RQ1: correspondence and approximation diagnostics}

At matched checkpoints and the same $\vomega$, we construct the following
scores:
\begin{enumerate}
    \item exact or fully unrolled hypergradient on the tractable problem;
    \item converged finite-penalty VF score, Eq.~\eqref{eq:vf-score};
    \item truncated/warm-start VF score after a controlled number of inner
    updates;
    \item ideal RHO score with an updating joint comparator;
    \item practical RHO score with a frozen holdout-only IL model.
\end{enumerate}
We report signed cosine similarity, Spearman rank correlation, Top-$B$
overlap, and selected-noise fraction. We sweep the penalty multiplier
$\alpha$, or a preregistered schedule $\alpha_r$, and report the associated
inner residuals. This separates finite-penalty bias from optimization and
frozen-comparator errors.

\subsection{RQ2: one-shot versus multi-round optimization}

All methods begin from the same model initialization and the same random
budget-$B$ subset $S_0$. After the common warm start, we compare:
\begin{center}
\begin{tabular}{lll}
\toprule
ID & Method & Purpose \\
\midrule
U-fixed & fixed random subset & no reselection baseline \\
U-round & random subset each round & reselection-only baseline \\
RHO-1 & one practical RHO selection & one-shot practical attribution \\
VF-1 & one verified VF step & one-shot value-function direction \\
RHO-M & repeated practical RHO & practical multi-round selection \\
VF-M & persistent projected VF & iterative outer optimization \\
\bottomrule
\end{tabular}
\end{center}
The outer update interval is the number of real target epochs between score
updates,
\begin{equation}
\tau\in\{1,5,10,20,\infty\},
\end{equation}
where $\tau=\infty$ performs one selection only after the prescribed target
training block. This variable replaces the earlier ``persistence horizon''
and ``online steps'' sweeps: it directly controls when a new selector state can
affect the real training trajectory. At each finite $\tau$, the continuous
state persists; only a separate reset ablation returns it to
$\vomega^{\mathrm{unif}}$.

The primary estimand is the paired multi-round gain
\begin{equation}
\Delta_{\mathrm{multi}}
=A(\mathrm{VF\mbox{-}M})-A(\mathrm{VF\mbox{-}1}),
\end{equation}
at equal target data and optimizer-step budgets. We also report target
accuracy versus selected examples, accuracy versus wall-clock/compute,
selected-noise fraction, class coverage, and selection turnover.

\subsection{RQ3: nested approximation ablations}

Starting from verified VF-M, we introduce one approximation at a time:
\begin{enumerate}
    \item reduce inner updates while keeping $\alpha$ and $\vomega$ fixed;
    \item freeze the joint comparator, then replace it with a holdout-only IL
    model;
    \item replace global scoring with RHO's batch-local $n_B\!\to n_b$
    protocol;
    \item delay Top-$B$ discretization while retaining continuous updates;
    \item reset the continuous state after each selection round.
\end{enumerate}
We additionally compare capped-simplex projection with the budget-calibrated
sigmoid parameterization from Section~\ref{sec:formulation}; this is a
parameterization ablation, not a different selection budget.
Four diagnostic controls isolate the multi-round mechanism: frozen target
state, frozen inner states, reset $\vomega$, and delayed Top-$B$. Their paired
differences decompose gains due to changes in the target trajectory, inner
tracking, persistent outer optimization, and repeated discretization.

Target and inner batch sizes are separate variables. We first screen
\texttt{target\_batch\_size}$\in\{32,128,\text{full selected set}\}$ with a
fixed inner batch, then screen
\texttt{inner\_batch\_size}$\in\{32,128,320,\text{full candidate set}\}$ at the
chosen target batch. Comparisons match examples processed and separately
report optimizer steps. Only the two best configurations are repeated with
all seeds.

The faithful practical RHO baseline follows the original online protocol:
sample a large batch of $n_B=320$, compute frozen-IL reducible losses, select
$n_b=32$, and immediately update the target on those examples. Selection and
target-update BatchNorm statistics are computed on the large and selected
batches, respectively. Global RHO scoring is reported as our extension rather
than as the original RHO-LOSS algorithm.

\subsection{Datasets, budgets, and staged execution}

We proceed from MNIST to CIFAR-10 and CIFAR-100. The initial study uses clean
and 10\% symmetric-label-noise conditions, 10\% retention, and three paired
seeds. We expand to 20\% and 30\% noise only after the direction is stable
across at least two of three seeds. After fixing the 10\% protocol, we evaluate
retention in $\{5\%,10\%,15\%,20\%\}$ for U-round, VF-1, VF-M, and the best
practical RHO variant.

Every outer update stores the continuous weights, selected indices, raw and
standardized scores, weight sum/minimum/maximum/effective sample size, both
inner objectives and residuals, score cosine, Top-$B$ Jaccard similarity,
class histogram, selected-noise fraction, target metrics, processed examples,
optimizer steps, wall-clock time, and peak memory. A run is admissible only
if its configuration, seeds, initial subset $S_0$, and all state transitions
are reproducible.
